\section{Grupo fundamental}

\subsection{Arcos y homotopias}
Para la construccion de nuestra invariante, debemos de familiarizarnos
con los arcos y las deformaciones continuas entre ellos llamadas
\emph{homotopias}.

\begin{definicion}[Arco]
  Un arco en el espacio \(X\) es una funcion continua \(f : I \to X \)
  donde \(I = [0,1]\)
\end{definicion}
La definicion tradicional es ligeramente mas general cambiando que \(I\)
sea un subconjunto convexo acotado de un espacio unidimensional, pero
siempre podemos recuperar esta definicion rescalando la distancia.

\begin{definicion}[Homotopia]
  Dados dos arcos \(f,g : I \to X\), diremos que \(f\) es homotopico a
  \(g\) si existe una funcion continua \(F : I \times I \to X \) tal que
  \[ \begin{matrix}
      F (x, 0) = f(x) & F (x, 1) = g(x)
     \end{matrix}
  \]
  Donde \(F\) sera llamada una homotopia entre \(f\) y \(g\). La
  existencia de dicha relacion entre dos funciones sera denotada por \(f
  \homRelAlt g\).
\end{definicion}
\begin{figure}[h]
  \centering
  \includegraphics[scale=0.5]{./imagenes/homotopia.png}
  \caption{caracterizacion homotopia}
  \label{fig:homotopia-entre-funciones}
\end{figure}
Podemos pensar en el segundo argumento de una homotopia como el grado de
deformación entre dos funciones.
Si tratamos con arcos \(f,g : I \to X\) que posean los mismos puntos
iniciales y finales, es decir \(f(0) = g(0) = x_0, \; f(1) = g(1) =
x_1 \) podemos definir una relacion homotopica ligeramente mas fuerte.

\begin{definicion}[Arco homotopia]
  \(f,g : I \to X\) son \emph{arco homotopicas} entre sí, si tienen los mismos
  puntos inicial y final \(x_0, x_1\) respectivamente y existe una homotopia entre
  ellos tal que cumpla
  \[
    \begin{matrix}
      F(s,0) = f(s) & F(s,1) = g(s) \\
      F(0,t) = x_0  & F(1,t) = x_1
    \end{matrix}
  \]
  La existencia de esta relacion entre dos funciones sera denotada por
  \(f \simeq g\).
\end{definicion}
% todo(slack) hacer algun dibujo de las homotopias entre las curvas
Una vez definida estas relaciones, una pregunta natural es si estas definen
una relacion de equivalencia, ya que nos gustaria identificar clases de
curvas como elementos de alguna teoria algebraica. La respuesta es
afirmativa para ambas relaciones (el punto de inicio y final no juegan
papel), se mostrara para las homotopias.

\begin{teorema}
  \(\homRelAlt\) es una relacion de equivalencia
\end{teorema}
\begin{proof}
  Hemos de probar que esta relacion cumple la reflexividad, simetria y
  transitividad. Sean \(f,g,h : I \to X\) tres arcos arbitrarios.
  \begin{itemize}
  \item La reflexividad es directa pues la funcion \(F(x,t) = f(x)\) es
    una deformacion continua de \(f\) a \(f\), por tanto \(f \homRelAlt
    f\).

  \item La simetria se obtiene de invertir el grado de deformacion de la
    homotopia original. Formalmente, dado \(f \stackrel{.}{\simeq} g\)
    tenemos una homotopia entre estas \((x,t) \mapsto F(x,t)\). A partir
    de aqui podemos definir
    \begin{equation}
      \label{eq:homotopy-simetry}
      (x,t) \mapsto \hat{F}(x,t) := F(x,1-t)
    \end{equation}
    con \(\hat{F}\) una homotopia entre \(g\) y \(f\). Por tanto \(g
    \stackrel{.}{\simeq} f\).


  \item La transitividad se obtiene a partir de de dividir \(I\) en dos
    intervalos donde se deformen individualmente cada homotopia al doble
    del grado. Formalmente dado \(f \stackrel{.}{\simeq} g\) y \(g
    \stackrel{.}{\simeq} h\) representadas por las homotopias \(F\) y
    \(G\) respectivamente, se define
    \[ FG(x,t) = \begin{cases}
        F(x,2t) & t \in [0,\frac{1}{2}] \\
        G(x,2t - 1) & t \in [ \frac{1}{2} , 1]
      \end{cases}
    \]
    Esta es una deformacion continua claramente en \((x,t) \in I \times
    [0, \frac{1}{2}) \cup I \times (\frac{1}{2}, 1]\). La continuidad en
    \(I \times \{\frac{1}{2}\}\) proviene de la consistencia en dicho
    punto de ambas homotopias
    \[ F(x,2 \cdot \frac{1}{2}) = g(x) = G(x, 2 \cdot \frac{1}{2} - 1)\]
    lo que nos permite utilizar el lema del pegamiento para obtener la
    continuidad de \(FG\). Obteniendo asi \(f \stackrel{.}{\simeq} h\).
  \end{itemize}
\end{proof}

En analisis real, es comun trabajar con conjuntos convexos, en arcos
sobre estos, veremos una homotopia usual repetidamente que es llamada
homotopia de linea recta.
\begin{definicion}\label{def:homotopia-linea}
  Sean \(f,g : I \to X\) dos arcos arbitrarios. Si \(f(I),g(I) \subset
  \mathcal C\), subconjunto convexo de \(X\), entonces
  \[ (x,t) \mapsto F(x,t) = (1-t) \cdot f(x) + t \cdot g(x) \]
  Es una homotopia entre \(f\) y \(g\), la cual es continua en virtud de
  ser combinacion lineal de funciones continuas.
\end{definicion}

Con estas definiciones ya podemos empezar a hablar de \([f]\) clases
de equivalencia de funciones bajo una relacion homotopica
\[ [f] = \{ g : I \to X \mid f \stackrel{.}{\simeq} g \} \]
Notar que la construccion es valida para las dos relaciones aunque por
nuestra meta de construir el grupo fundamental nos interesa
principalmente la relacion arco homotopica, por razones a estudiar mas
adelante.

\subsection{Grupoide de arcos}
Si tenemos dos arcos continuos \(f,g\) tales que el punto final de \(f\)
sea el punto inicial de \(g\) nos da la idea de que podemos construir un
camino continuo que recorra \(f\) luego recorriendo \(g\). Esta idea es
conocida formalmente como el producto de dos arcos.

\begin{definicion}[Producto de arcos]
Para dos arcos \(f,g : I \to X\) tales que
\(f(1) = g(0)\), se define el producto \(f * g \)
\[ (f*g) (s) = \begin{cases}
    f(2s) & s \in [0,\frac{1}{2}] \\
    g(2s - 1) & s \in [\frac{1}{2} , 1]
  \end{cases}
\]
la cual sigue siendo una funcion continua en virtud del lema del
pegamiento.
\end{definicion}

Esta construccion se puede reutilizar para clases
\emph{arco}-homotopicas \([f],[g]\) que compartan punto final e inicial
respectivamente para definir un producto de clases de equivalencia
\[ [f] * [g] := [f * g]\]
El cual esta bien definido pues si \(f \simeq_{ah} f'\), \(g \simeq_{ah}
g'\) a traves de \(F, G\) respectivamente
\[H(s,t) = \begin{cases}
    F(2s,t) & s \in [0, \frac{1}{2}] \\
    G(2s - 1, t) & s \in [\frac{1}{2} , 1]
  \end{cases}
\]
Es la homotopia que relaciona a los arcos que sean homotopicos a
\(f*g\). Esta ademas es continua en virtud otra vez del lema del pegamiento.

Con esta operacion binaria, una pregunta a hacerse es si tiene
\((X/\simeq_{ah} , *)\) posee
estructura de grupoide. Esto equivale a cumplir 3 propiedades
\begin{enumerate}
\item \textbf{Asociatividad} Si \([f] * ([g] * [h])\) esta definido entonces
  tambien lo debe estar \(([f] * [g]) * [h]\) y deben de ser iguales.
\item \textbf{Identidades izquierda y derecha} Para todo \([f]\) con
  \(x_0, x_1\) puntos inicial y final respectivamente, debe de
existir elementos \([k_{0}], [k_{1}]\) tales que
\[ \begin{matrix}
    [f] * [k_{1}] = [f] & & [k_{0}] * [f] = [f]
  \end{matrix}
\]
\item \textbf{Inverso} Para todo \([f]\) clase de arcos con \(x_0, x_1\)
  puntos inicial y final respectivamente debe de existir un elemento
  \([\bar{f}]\) que cumpla
\[ \begin{matrix}
    [f] * [\bar{f}] = [k_{x_0}] & & [\bar{f}] * [f] = [k_{x_1}]
  \end{matrix}
\]
\end{enumerate}
Para probar esto necesitamos primero definir a nuestros candidatos de
\(k_{x_1}, k_{x_2}, \bar{f}\) ademas de algunas funciones auxiliares,
iniciando por los mapeos constantes e identidad en \(I\)
\[ \begin{matrix}
     e_0 : & I \to I & e_0(t) := 0 \\
     e_1 : & I \to I & e_1(t) := 1 \\
     i :   & I \to I & i(t) := t \\
     \bar{i} : & I \to I & i(t) := 1 - t
   \end{matrix}
   \]
Para todo arco \(f : I \to X \), el elemento \(\bar{f} : I \to X \) esta
definido (en el espiritu de \eqref{eq:homotopy-simetry}) por
\[ \bar{f} (s) := f (1 - s) \]
Para todo \(x \in X \) se define la curva constante
\begin{align*}
  k_x : &I \to X \\
        &t \mapsto x
\end{align*}
Con esto podemos afrontar la demostración
\begin{proof}
Se procedera \(2 \to 3 \to 1\) con lemas en medio. Sea \(f : I \to X\) el
representante de \([f]\). sean \(x_0, x_1\) los puntos inicial y final de
\(f\) respectivamente.

\paragraph{(2).} Dados \(i,e_1\) arcos en \(I\) definidos anteriormente,
es claro que \(i * e_1\) es tambien un arco (continuo) en \(I\); mas aun,
ya que \(I\) es convexo se tiene \(i * e_1 \simeq_{ah} i\) con la
homotopia de la linea recta entre estos. Se desprende de esto ultimo que
\( [i] = [i * e_1]\). Para continuar necesitamos el siguiente lema
\begin{lema}[Distributividad de la composicion sobre producto.]
\label{lema:dist-composicion-producto}
\[\forall a,b : I \to I,\ f \circ (a * b) = (f \circ a) * (f \circ b) \]
\end{lema}
\begin{proof}
  \[ f \circ (a*b) (s) =
    \begin{cases}
      f (a(2s)) & s \in [0,\frac{1}{2}] \\
      f (b(2s - 1)) & s \in [\frac{1}{2} , 1]
    \end{cases}
    = (f \circ a) * (f \circ b)
  \]
\end{proof}
Recordemos ademas que la composición de funciones continuas es continua.
Luego esto nos permite afirmar en especifico que
\[ f \circ (i * e_1) \simeq_{ah} (f \circ i) * (f \circ e_1) \] en virtud
de la composición de \(f\) con la homotopia de linea recta entre \(i *
e_1 \) y \(i\). Obteniendo asi que
\[ [f \circ (i * e_1)] = [(f \circ i) * (f \circ e_1)] \] pero al reducir
las composición obtenemos
\[
  \begin{matrix}
    f \circ i = f & f \circ e_1 = k_{x_1} & f \circ (i * e_1) = f *
    k_{x_{1}} \\
  \end{matrix}
\]
\[ \implies [f] * [k_{x_1}] := [f * k_{x_1}] = [f \circ (i * e_1)] = [ f
  \circ i] = [f] \]
probando (2).

\paragraph{(3).} De manera similar a la pregunta anterior, notar que \(i,
\bar{i}\) definidos anteriormente cumplen
\[ i * \bar{i} \simeq_{ah} e_0 \]
por la homotopia de linea recta en el espacio \(I\) que es convexo.
Utilizando el lemma \eqref{lema:dist-composicion-producto} tenemos que
\[ f * \bar{f} = f \circ (i * \bar{i}) \simeq_{ah} f \circ e_0 =
  k_{x_0} \]
\[ \implies [f] * [\bar{f}] = [k_{x_0}] \]
De manera analoga se prueba la inversa por la izquierda.

\paragraph{(1).}
% todo(slack)
\end{proof}


\subsection{Grupo fundamental}
Sobre nuestros arcos tenemos una estructura de grupoide porque no siempre
podemos hacer coincidir los puntos inicial y final. Si nos fijaramos en
un solo punto \(x_0 \in X\) tal que todas los arcos \(f : I \to X\)
cumplieran que
\[ f(0) = x_0 = f (1) \]
Entonces las propiedades que demostramos anteriormente nos dirian que el
espacio \((X / \simeq_{ah}, *)\) es un grupo.
\begin{definicion}[Grupo fundamental]
  Sea \(X\) un espacio topologico y \(x_0 \in X\) un punto. La coleccion
  \[ \pi_1 (X,x_0) = \{ [f] \mid f : I \to X, f \text{ continua}\}\]
  \[ [f] := \{ g : I \to X \mid f \simeq_{ah} g\}\]
  junto con el producto de caminos \((*)\) forman el llamado
  \textbf{grupo fundamental} \((\pi_1(X,x_0), (*))\)
\end{definicion}
Nuestro interes es mostrar que para un espacio dado, su grupo fundamental
es una invariante topologica. Esto quiere decir que dado dos espacios puntuados
\((X,x_0)\) e \((Y,y_0)\) homeomorfos, es decir que existe un isomorfismos continuo
entre ellos, debe de existir un respectivo isomorfismo de grupos entre
\(\pi_1(X,x_0)\) y \(\pi_1(Y,y_0)\). Para esto debemos recordar que
significa ser isomorfo entre grupos
\begin{definicion}
  Sean \((G,(*)), (G', (\cdot))\) does grupos. \(F : G \to G'\) es un
  homomorfismo entre grupos si y solo si
  \[ \forall x,y \in G,\ F (x * y) = F(x) \cdot F(y)\]
  \[ e \text{ identidad en } G \implies F(e) \text { identidad en } G' \]
\end{definicion}

\paragraph{Ejemplo \(\mathbb{R}^2 - \{(1,1)\}\).} Este espacio es
relativamente sencillo de visualizar, haremos un argumento intuitivo para
ver su grupo fundamental. Sea \(f : I \to \mathbb{R}^2 - \{(1,1)\}\) un
camino arbitrario

%% old
Una pregunta natural ha
hacerse es si la eleccion de punto influye en el tipo de grupo que se
obtiene en \((X / \simeq_{ah}, *)\). En general la respuesta es
afirmativa, lo cual complica nuestra tarea puesto que queremos definir
una invariante topologica que dependa del espacio total. Pero podemos
recuperar nuestro objetivo si nos restringimos a espacios \(X\) que sean
arco-conexos.
\begin{definicion}[Arco-conexidad]
  Diremos que el espacio \(X\) es \emph{arco conexo} si para todo par de
  puntos \(x,y \in X\) existe un arco continuo \(f : I \to X\) tal que
  cumpla
  \[ f(0) = x \qquad f(1) = y \]
\end{definicion}
Para estos espacios, denotamos a \((X / \simeq_{ah}, *)\) como
\((\pi_1(X,x_0), *)\) y le llamaremos \textbf{grupo fundamental} de \(X\) con
respecto a \(x_0\). Es claro que la importancia del punto \(x_0\) es poca
considerando la arco-conexidad. Puesto que si tengo \([f] \in
\pi_1(X,x_0)\) y \(\alpha : I \to X\) un arco entre \(x_0\) y \(x_1\),
existe un camino \([\hat{\alpha} * f * \alpha] \in \pi_1(X,x_1)\), es
decir un isomorfismo entre \(\pi_1 (X, x_0)\) y \(\pi_1 (X, x_1)\). Aun
asi es conveniente reificar los puntos bases de los grupo puesto que no
necesariamente existe un camino canonico entre ellos.
% La razon de porque aun guardamos a \(x_0\) como
% informacion del grupo es que a priori los isomorfismos entre \(\pi_1 (X,
% x_0)\) y \(\pi_1 (X, x_1)\) pueden depender del arco \(\alpha\) elegido,
% de los cuales ningun

%todo(slack): ejemplo de R^2 sin el punto
%todo(slack): ver la discucion de la propiedades functoriales (talvez?)